\documentclass[11pt]{article}
\usepackage[margin=1in]{geometry}
\usepackage{amsmath,amssymb,amsthm,hyperref}
\newcommand{\E}{\mathbb E}
\newcommand{\Prob}{\mathbb P}
\newcommand{\T}{\mathbb T}
\title{HOCF--RIESZ: Round 003\\The initial iid pair and collision interface}
\author{Campaign research memorandum}
\date{17 September 2026}
\begin{document}
\maketitle
\noindent Published input checkpoint:
\texttt{2f532ccd2cb5f3d84db456be96e438f03f7e2ad2}.
The initial pair constructions used isolated worktrees based on the
published Round 001 input; later proofs and reviews used the Round 002
checkpoint above. This memorandum records initial iid statements. The round
report gives the actual per-claim audit status. Neither the bare potential
nor its truncated version is identified with the backward corrector, and
no positive-time interacting-law estimate is inferred.

\section{An exact finite-particle projection formula}
Let $N\ge2$, let $X_1,\ldots,X_N$ be iid with probability law $\mu$, and let
$\Phi$ be real, symmetric and in $L^2(\mu\otimes\mu)$. Set
\[
 \theta=\mu^2(\Phi),\qquad
 h(x)=\int\Phi(x,y)\,d\mu(y)-\theta,\qquad
 H(x,y)=\Phi(x,y)-\theta-h(x)-h(y).
\]
The canonical kernel $H$ has conditional mean zero in either variable.
Orthogonality gives
\[
 \|\Phi\|_2^2=\theta^2+2\|h\|_2^2+\|H\|_2^2.
\]
With the campaign's ordered distinct-label denominator $N^2$, define
\[
 P_N[\Phi]=\frac1{2N^2}\sum_{i\ne j}\Phi(X_i,X_j)
       -\frac1N\sum_i\int\Phi(X_i,y)\,d\mu(y)+\frac\theta2.
\]
Collecting constants and single-label terms gives the exact identity
\begin{equation}\label{projection}
 P_N[\Phi]=-\frac\theta{2N}
      -\frac1{N^2}\sum_i h(X_i)
      +\frac1{N^2}\sum_{i<j}H(X_i,X_j).
\end{equation}
The first projection remains because each label occurs only $N-1$ times
in the unordered pair sum. Conditional centering makes distinct canonical
pair terms orthogonal, including pairs sharing one label; it also makes
them orthogonal to the linear sum. Therefore
\begin{align}
 \E P_N[\Phi]&=-\frac\theta{2N},\label{mean}\\
 \E|P_N[\Phi]|^2
 &=\frac{\theta^2}{4N^2}+\frac{\|h\|_2^2}{N^3}
       +\frac{N-1}{2N^3}\|H\|_2^2\label{moment}\\
 &\le\frac{N-1}{2N^3}\|\Phi\|_2^2.\label{bound}
\end{align}
The last constant is attained by a nonzero canonical kernel whenever the
probability space admits one; at $N=2$ equality holds for every kernel.
A constant kernel yields a deterministic $-\theta/(2N)$, not zero.
A kernel $h(x)+h(y)$ with centered $h$ yields precisely the linear term in
\eqref{projection}. These test the centering and the denominator.

There is no self-label evaluation $\Phi(X_i,X_i)$. For atomless $\mu$,
spatial-diagonal values may be changed freely. With atoms, distinct labels
can coincide spatially, and diagonal values at those atoms belong to the
$L^2(\mu^2)$ equivalence class. Deleting labels does not delete these values.

\section{Riesz normalization and the square-integrability boundary}
Now use Haar probability measure on the unit torus and $0<s<d$. Write
\[
 c_{d,s}=\pi^{s-d/2}\frac{\Gamma((d-s)/2)}{\Gamma(s/2)},\qquad
 \widehat g(k)=c_{d,s}|k|^{s-d}\quad(k\ne0),\qquad \widehat g(0)=0.
\]
Let $p_t$ be the torus heat kernel with multiplier
$e^{-4\pi^2t|k|^2}$. Its Gaussian periodization gives the $L^1$ representation
\begin{equation}\label{heat}
 g(x)=\frac{4^{(d-s)/2}\pi^{d/2}}{\Gamma(s/2)}
       \int_0^\infty t^{(d-s)/2-1}(p_t(x)-1)\,dt.
\end{equation}
At small $t$ the $L^1$ norm of $p_t-1$ is at most two, and at large $t$ it
decays exponentially. Fubini computes exactly the prescribed Fourier
coefficient. The whole-space Gaussian term in \eqref{heat} integrates to
$|x|^{-s}$ with coefficient one. Nonzero lattice translates give integrable
$e^{-c/t}$ bounds near $t=0$, with all derivatives; the large-time difference
is also smooth. Hence locally
\[
 g(x)=|x|^{-s}+H_0(x),\qquad H_0\text{ smooth near }0.
\]
In particular $g$ is integrable and bounded below, and belongs to $L^2$
exactly when $2s<d$.

For Haar iid samples the pair kernel $g(x-y)$ is canonical in $L^1$, so
\[
 P_N[g]=N^{-2}\sum_{i<j}g(X_i-X_j),\qquad \E P_N[g]=0.
\]
It is finite almost surely for every finite $N$. If $2s\ge d$, its second
moment is infinite. Indeed, a lower bound $g\ge-C$ implies, with
$M=N(N-1)/2$,
\[
 (P_N[g])_+\ge N^{-2}\big(g(X_1-X_2)-C(M-1)\big)_+,
\]
whose square has infinite expectation. This argument does not expand
undefined cross moments and does not assert nonconvergence in probability.

\section{Heat cutoffs and two different spatial scales}
Let $g_\varepsilon=p_\varepsilon*g$ and
$b_N=\min(\beta_N,1)$, $\sigma_N=\sqrt{Nb_N}$. Parseval and
\eqref{moment} give
\begin{align}
 \|g_\varepsilon\|_2^2
 &=c_{d,s}^2\sum_{k\ne0}|k|^{2s-2d}
                       e^{-8\pi^2\varepsilon|k|^2},\label{sum}\\
 \E|\sigma_NP_N[g_\varepsilon]|^2
 &=\frac{b_N(N-1)}{2N^2}\|g_\varepsilon\|_2^2.\label{scaled}
\end{align}
Counting lattice points in max-norm shells in \eqref{sum} gives, with
positive constants depending on $d,s$ and a fixed small cutoff range,
\[
 \|g_\varepsilon\|_2^2\asymp
 \begin{cases}
  1,&2s<d,\\
  \log(1/\varepsilon),&2s=d,\\
  \varepsilon^{-(2s-d)/2},&2s>d.
 \end{cases}
\]
The shell $|k|_\infty=n$ has $(2n+1)^d-(2n-1)^d$ points, between
$2dn^{d-1}$ and $2d3^{d-1}n^{d-1}$. On this shell
$n\le|k|\le\sqrt d n$, so the sum is bounded above and below by explicit
constants times $\sum_{n\ge1}n^{2s-d-1}e^{-a\varepsilon n^2}$,
with $a=8\pi^2$ and $8\pi^2d$, respectively. Integral comparison gives
the displayed three orders. The full proof in the sealed pair memorandum
gives explicit two-sided constants and their cutoff interval.

At $\varepsilon_N=N^{-2/d}$, the high-singularity scaled second moment is
of order $b_NN^{2s/d-2}$ and tends to zero for every $s<d$. At
$\varepsilon_N=N^{-4/d}$ it is of order $b_NN^{4s/d-3}$ when $2s>d$.
These cutoffs correspond respectively to lengths $N^{-1/d}$ and
$N^{-2/d}$. A regularized variance tending to zero does not control the
unregularized statistic: for $2s\ge d$ the latter and its difference from
every fixed smooth cutoff have infinite second moment. A probability
comparison needs a separate estimate.

\section{A probability estimate that retains close pairs and the mean}
Fix $0<r<1/4$. Let
\[
 h_r=g\mathbf1_{\{\operatorname{dist}(x,0)\ge r\}},\qquad
 \tau_r=\int_{\operatorname{dist}(x,0)<r}g(x)\,dx,\qquad
 k_r=h_r+\tau_r.
\]
Then $k_r$ is canonical. Let $E_r$ be the event that no pair of distinct
labels is within distance $r$, and let $v_d$ be the unit-ball volume.
A union bound and the exact normalization give
\[
 \Prob(E_r^c)\le\frac{N(N-1)}2v_dr^d,
 \qquad
 P_N[g]=P_N[k_r]-\frac{N-1}{2N}\tau_r\quad\text{on }E_r.
\]
The mean correction has coefficient $(N-1)/(2N)$, not zero. Put
\[
 V_r=\|k_r\|_2^2,\qquad
 B_N(r)=\sigma_N\frac{N-1}{2N}|\tau_r|.
\]
For $z>B_N(r)$, Chebyshev and \eqref{scaled} for this canonical kernel yield
\begin{equation}\label{prob}
 \Prob\{|\sigma_NP_N[g]|>z\}
 \le\frac{N(N-1)}2v_dr^d+
       \frac{b_N(N-1)V_r}{2N^2(z-B_N(r))^2}.
\end{equation}
The local singularity supplies
\[
 \tau_r=\frac{dv_d}{d-s}r^{d-s}+O(r^d),\qquad
 V_r\le C_{d,s}
 \begin{cases}
 1,&2s<d,\\
 1+\log(1/r),&2s=d,\\
 r^{d-2s},&2s>d.
 \end{cases}
\]
The constants may equivalently be expressed by the displayed smooth
remainder and the integral outside the fixed ball; they are independent
of $N,\beta_N,r$.

At $2s=d$, choose $r_N=N^{-3/d}$. For $b_N=1$, close-pair probability,
scaled discarded mean, and truncated variance are respectively
$O(N^{-1})$, $O(N^{-1})$, and $O(\log N/N)$. For $d/2<s<3d/4$, choose
\[
 r_N=N^{-a},\qquad \frac2d<a<\frac1{2s-d}.
\]
All three errors in \eqref{prob} then vanish. Together with the direct
$L^2$ result below $d/2$, this proves
\[
 0<s<\frac{3d}4\quad\Longrightarrow\quad
 \sqrt N P_N[g]\longrightarrow0\quad\text{in probability}.
\]
It also holds at $\sigma_N\le\sqrt N$ for every positive temperature
sequence in this range.

For $d/2<s<d$ a further sufficient condition is
\[
 b_NN^{4s/d-3}\longrightarrow0.
\]
Indeed put $A_N=\sqrt{b_N}N^{2s/d-3/2}$ and choose
$r_N=N^{-2/d}A_N^{1/(2s-d)}$. The three controlling factors are
\[
 N^2r_N^d=A_N^{d/(2s-d)},\quad
 \sigma_Nr_N^{d-s}=A_N^{s/(2s-d)},\quad
 \frac{b_N}{N}r_N^{d-2s}=A_N.
\]
This proves the asserted probability convergence. It supplies no converse.
At $s\ge3d/4$ and $b_N=1$, requiring no close pairs forces
$r_N=N^{-2/d}\ell_N$ with $\ell_N\to0$, whereas the scaled truncated
variance grows as $N^{4s/d-3}\ell_N^{-(2s-d)}$. This disproves closure of
that particular truncation strategy, not convergence of the original
statistic. An endpoint or stable-law claim would require a new proof.

\section{An absolute-moment estimate and a sharp converse}
Center both parts of the spatial split: $k_r=g\mathbf1_{\{\rho\ge r\}}+
\tau_r$ and $j_r=g\mathbf1_{\{\rho<r\}}-\tau_r$, where $\rho$ denotes
distance to zero. Then $g=k_r+j_r$. Put
$D(r)=\int_{\rho<r}|g|\le C_1r^{d-s}$. The exact outer variance and the
triangle inequality over the unordered inner sum give
\[
 \E|\sigma_NP_N[g]|
 \le \sqrt{\frac{b_N}{2N}}\|g\mathbf1_{\{\rho\ge r\}}\|_2
       +\sqrt{Nb_N}\,D(r).
\]
Indeed $\|j_r\|_1\le2D(r)$ and
$\E|P_N[j_r]|\le (N-1)\|j_r\|_1/(2N)$. Choose $r=N^{-2/d}$.
For fixed $d,s$ this proves, uniformly in every positive temperature,
\begin{equation}\label{L1rates}
 \E|\sigma_NP_N[g]|\le C_{d,s}
 \begin{cases}
  \sqrt{b_N/N},&2s<d,\\
  \sqrt{b_N(1+\log N)/N},&2s=d,\\
  \sqrt{b_N}\,N^{2s/d-3/2},&2s>d.
 \end{cases}
\end{equation}
The first row follows directly from the untruncated second moment.
No exclusion of close pairs is used in this argument. It was constructed
separately in the statement-only probability reconstruction.

Here is the new converse for $d/2<s<d$. Let $u_N=N^{2s/d-2}$ and
$R_N=N^{-2/d}$. Fix a sufficiently small radius $R$ such that
$g(x)\ge\tfrac12|x|^{-s}$ for $0<|x|<R$,
$\tau_r\le C_\tau r^{d-s}$ and $V_r\le C_Vr^{d-2s}$ for $r<R$.
For large $N$ take $R_N<R$ and count pairs at distance less than
$\delta R_N$, with a fixed $0<\delta<1$:
\[
 Z_N=\sum_{i<j}\mathbf1_{\{|X_i-X_j|<\delta R_N\}},\qquad
 m_N=\E Z_N=\frac{N(N-1)}2v_d\delta^dN^{-2}.
\]
Distinct pair indicators are pairwise independent: for shared labels,
condition on that label and use Haar translation invariance. Joint
independence of all the indicators is neither true in general nor used.
Thus $\E Z_N^2\le m_N+m_N^2$, and Cauchy--Schwarz yields
\[
 \Prob(Z_N>0)\ge\frac{m_N}{1+m_N}\ge c_0\delta^d,
 \qquad c_0=\frac{v_d}{4(1+v_d/2)}.
\]
The exact outer--inner decomposition is
\[
 P_N[g]=P_N[k_{R_N}]+N^{-2}\sum_{i<j}
    g(X_i-X_j)\mathbf1_{\{|X_i-X_j|<R_N\}}
       -\frac{N-1}{2N}\tau_{R_N}.
\]
All inner summands are positive. On $Z_N>0$ at least one contributes
$\tfrac12\delta^{-s}u_N$. If $\delta^{-s}\ge2C_\tau$, the last two
terms are at least $\tfrac14\delta^{-s}u_N$. Meanwhile
\[
 \E P_N[k_{R_N}]^2\le\frac{C_V}{2}u_N^2,
 \quad
 \Prob\left(P_N[k_{R_N}]<-\tfrac18\delta^{-s}u_N\right)
       \le32C_V\delta^{2s}.
\]
Since $2s>d$, choose the same fixed $\delta$ so that
$\delta^{2s-d}\le c_0/(64C_V)$. Subtracting the last failure
probability from the close-pair probability, without assuming event
independence, proves
\begin{equation}\label{sharp}
 \Prob(P_N[g]\ge a u_N)\ge p_0,
 \qquad a=\delta^{-s}/8,\quad p_0=c_0\delta^d/2.
\end{equation}
Consequently $ap_0u_N\le\E|P_N[g]|\le C_{d,s}u_N$. Multiplying
by $\sigma_N$ shows that, in this range, vanishing in probability is
equivalent to vanishing in $L^1$, and equivalent to
\[
 B_N:=b_NN^{4s/d-3}\longrightarrow0.
\]
For necessity, a subsequence on which $B_N\ge\eta>0$ retains the fixed
positive tail probability in \eqref{sharp}. Bounded $B_N$ gives tightness
by the absolute-moment upper bound. From an unbounded scalar subsequence
extract a further one tending to infinity; the same lower probability
then proves non-tightness. At $\sqrt N$ scale the bare statistic is
therefore negligible below $3d/4$, tight and nonvanishing at $3d/4$, and
not tight above it. No endpoint distribution or stable law is identified.
This converse is a separate assertion from the earlier sufficient-bound
argument; the round report records its actual independent review status.

\section{An exact model retaining internal pair transport}
The raw-potential threshold cannot automatically be applied to the
backward corrector. A solvable model makes the distinction concrete.
Let $A$ be a fixed symmetric matrix, $z=r\theta\ne0$,
$p=s+2$ and $\tau=T-t$. Retain only the Euclidean internal pair drift:
\[
 \left(\partial_t+\frac{2s}{N}r^{-s-2}z\cdot\nabla_z\right)\Phi_N
       =-s r^{-s}\theta^TA\theta,\qquad \Phi_N(T,z)=0.
\]
The factor two comes from $\nabla_x-\nabla_y=2\nabla_z$.
Directions are constant along characteristics and
$r(h)^p=r^p+2sp h/N$. Integrating the source along the forward
characteristic gives the unique solution in the absolutely continuous
characteristic class, with no equation or boundary condition imposed at
zero:
\begin{equation}\label{toy}
 \Phi_N(t,z)=\frac N4(\theta^TA\theta)
       \left[(r^p+2sp\tau/N)^{2/p}-r^2\right].
\end{equation}
Differentiation gives $-\partial_\tau F+(2s/N)r^{-s-1}\partial_rF
=-sr^{-s}$ for the radial amplitude $F$, checking both sign and coefficient.

Set $\ell=(2sp\tau/N)^{1/p}$ for $\tau>0$. Then
$F=(N/4)\ell^2 H_s(r/\ell)$, where
\[
 H_s(u)=(1+u^p)^{2/p}-u^2
 =\frac2p\int_0^1(u^p+v)^{-s/p}\,dv.
\]
This decreases from one. It is bounded above and below by positive
constants on $[0,1]$, and by constants times $u^{-s}$ on $[1,\infty)$.
The exact angular factor is
\[
 Q(A)=\int_{S^{d-1}}(\theta^TA\theta)^2\,d\theta
 =\frac{|S^{d-1}|}{d(d+2)}
       \big[(\operatorname{tr}A)^2+2\operatorname{tr}(A^2)\big].
\]
Thus every nonzero traceless symmetric $A$ also has a positive square norm.
For any fixed ball,
\[
 \|\Phi_N(t)\|_{L^2(B_R)}^2
  =\frac{Q(A)N^2\ell^{d+4}}{16}
       \int_0^{R/\ell}H_s(u)^2u^{d-1}\,du.
\]
When $\ell\le R$, the three orders are respectively
\[
 Q(A)\begin{cases}
  \tau^2R^{d-2s},&2s<d,\\
  \tau^2[1+\log(R/\ell)],&2s=d,\\
  N^{(2s-d)/(s+2)}\tau^{(d+4)/(s+2)},&2s>d.
 \end{cases}
\]
The full sealed proof supplies two-sided constants, the $\ell>R$ case,
and the zero-time endpoint. Despite finite radial amplitude, a nonscalar
$A$ has direction-dependent limits at collision. No smooth extension is
available in that case.

Choose a fixed smooth radial cutoff supported inside a torus ball of
radius less than $1/2$ and define the periodic pair diagnostic
$\mathcal H_{N,\tau}(x,y)=\chi(x-y)\Phi_N(T-\tau,x-y)$.
It is an $L^2$ kernel; its diagonal value is immaterial for a density.
For iid preparation with density bounded by $M$,
\[
 \|\mathcal H_{N,\tau}\|_{L^2(\mu^2)}^2
 \le M\|\chi\Phi_N\|_2^2.
\]
Equation \eqref{bound} now retains all centering contributions, and yields
\[
 \sup_{0\le\tau\le T}\E|\sigma_NP_N[\mathcal H_{N,\tau}]|^2
 \le C_T b_N\begin{cases}
  N^{-1},&2s<d,\\
  (1+\log N)/N,&2s=d,\\
  N^{-(d+2-s)/(s+2)},&2s>d.
 \end{cases}
\]
Every row vanishes for $s<d$. The supremum is outside the expectation.
This is an initial iid diagnostic for a zero-diffusion equation; positive
temperature appears only in the statistic's scale.

\section{What remains load-bearing}
These estimates hold at initial iid Haar preparation for the raw potential.
The probability result beyond the $L^2$ range illustrates why variance
divergence is not a probability obstruction. It does not identify the
backward pair kernel, compare singular and regularized particle dynamics,
or control any positive-time martingale. The sharp converse concerns only
the raw potential. The transport diagnostic, also available for bounded
densities, demonstrates a different initial endpoint power and does not
solve the full pair equation.

For the genuine pair corrector $\Psi_N$, a sufficient remaining initial
comparison is
\[
 \sqrt{b_N/N}\,
 \|\Psi_N(0)-\mathcal H_{N,T}\|_{L^2(\mu_0^2)}\longrightarrow0.
\]
This is open. The error equation must retain $2\nu\Delta_z$, ordinary
transport, both nonlocal response terms, the actual test/source and
periodic force remainder, and the cutoff term
\[
 (2s/N)r^{-s-1}\chi'(r)\Phi_N.
\]
For a nonscalar $A$, the leading
angular Laplacian has a term proportional to $r^{-2}$, so diffusion
cannot be added by an unproved small-coefficient argument. Well-posedness,
cutoff passage, evolved-law residuals and martingale estimates also remain
load-bearing. A weaker topology could replace the displayed comparison
only with its own proved fluctuation-scale estimate.

The separate transport review makes one limitation precise: if the
traceless part of $A$ is nonzero and $\tau>0$, the squared punctured
Laplacian contains a positive multiple of
$\int_0^\epsilon r^{d-5}\,dr$. Thus it is not locally $L^2$ in
dimensions $2,3,4$; in dimension two its absolute $L^1$ norm also diverges.
This disproves a proposed separate-norm perturbation estimate for this
profile, not existence of a different diffusion-retaining solution.

There is also an exact punctured rescaling of the local equation with
diffusion retained. Set $y=N^{1/(s+2)}z$ and
$\Phi_N=N^{s/(s+2)}F_N(\tau,y)$. The chain rule gives
\[
 \partial_\tau F_N=2\chi_N\Delta_yF_N+
       2s|y|^{-s-2}y\cdot\nabla_yF_N+s|y|^{-s}a(y/|y|),
 \qquad \chi_N=\frac{N^{2/(s+2)}}{\beta_N}.
\]
At critical microscopic coupling the exact identity is
\[
 \chi_N=\lambda_N^{-1}
       N^{s(s+2-d)/(d(s+2))}.
\]
Consequently this coefficient tends to zero, a positive constant, or
infinity according as $d>s+2$, $d=s+2$, or $d<s+2$.
This is only an operator rescaling on the punctured domain. No collision
domain, solution convergence, or limit exchange follows from the
coefficient limit. A future local comparison must handle its degree-zero
and degree-two angular modes in the appropriate parameter range.

The three regimes remain distinct:
\[
 \beta_NN^{2s/d-1}\to0,\qquad
 \lambda_N=\beta_NN^{s/d-1}\to0,\qquad
 \lambda_N\to\lambda\in(0,\infty).
\]
They are respectively the old energy-floor condition, full microscopic
subcriticality, and microscopic criticality. None is replaced by
the sufficient initial-pair condition above. The logarithmic interaction
remains outside this positive-power calculation.
\end{document}
